\documentclass[11pt]{article}
\usepackage[a4paper,margin=0.92in]{geometry}
\usepackage[T1]{fontenc}
\usepackage{lmodern}
\usepackage{amsmath,amssymb,amsthm,mathtools}
\usepackage{graphicx}
\usepackage{microtype}
\usepackage[hidelinks]{hyperref}
\usepackage{xcolor}
\usepackage{fancyhdr}
\setlength{\emergencystretch}{2em}

\definecolor{accent}{RGB}{48,76,101}
\newtheorem{theorem}{Theorem}
\newtheorem{lemma}[theorem]{Lemma}
\newtheorem{example}[theorem]{Example}
\newcommand{\T}{\mathbb T}
\newcommand{\D}{\mathbb D}
\newcommand{\A}{\mathcal A}
\newcommand{\dist}{\operatorname{dist}}
\newcommand{\m}{m}
\pagestyle{fancy}
\fancyhf{}
\setlength{\headheight}{14pt}
\lhead{Explicit Cauchy transforms under log--log entropy}
\rhead{Lane 4}
\cfoot{\thepage}
\fancypagestyle{plain}{%
  \fancyhf{}%
  \lhead{Explicit Cauchy transforms under log--log entropy}%
  \rhead{Lane 4}%
  \cfoot{\thepage}%
}

\title{\color{accent}Explicit Disk-Algebra Cauchy Transforms\\
under Log--Log Entropy\\[3pt]
\large A substantial partial answer to Limani--Malman}
\author{Automated Mathematics Run \texttt{fa\_banach\_001}, lane 4}
\date{27 August 2026}

\begin{document}
\maketitle

\begin{center}
\fbox{\parbox{0.91\linewidth}{\textbf{Status.}
\texttt{candidate\_substantial\_partial\_likely\_valid}.  An explicit
formula is obtained for a class strictly larger than the
Beurling--Carleson class.  The problem for an arbitrary closed set remains
open.}}
\end{center}

\section{Source problem}

In the Introduction of arXiv:2111.14112v3, PDF page~6, Limani and Malman
ask for an explicit formula for a nonzero measurable function $h$ supported
on a given closed set $E\subset\T$ such that
\[
 C_h(z)=\int_E\frac{h(\zeta)}{1-z\overline\zeta}\,d\m(\zeta)
\]
belongs to the disk algebra $\A$ \cite{LM}.  Their construction gives the
stronger conclusion $C_h\in\A^\infty$ when $E$ is a positive-measure
Beurling--Carleson set.

\begin{figure}[ht]
  \centering
  \includegraphics[width=\linewidth]{figures/open_problem_crop.png}
  \caption{The open statement in the Introduction of the source paper,
  PDF page~6.}
\end{figure}

We assume throughout that $E$ is proper and has positive normalized Lebesgue
measure; these are necessary for a nonzero integrable density supported on
$E$.  The case $E=\T$ is trivial.

\section{Statement and proof intuition}

Let $\{I_n\}$ be the component arcs of $\T\setminus E$, and set
\[
 \Phi(t)=\log\!\bigl(e+\log(1/t)\bigr),\qquad 0<t\leq1.
\]
The source construction places poles just outside Whitney subarcs and weights
a subarc of length $\ell$ by $\ell\log(1/\ell)$.  This pays enough entropy to
make the cutoff flat to every order.  For the disk algebra, flatness is not
needed.  We instead use the smaller weight $\ell\Phi(\ell)$.  The cutoff then
decays like a chosen negative power of $\log(e/\dist(\cdot,E))$.

On the complement of $E$, conjugate the cutoff and extend it by zero across
$E$.  A two-scale estimate shows that the extension has modulus of continuity
\[
 \omega(\delta)\lesssim
  \bigl(\log(e/\delta)\bigr)^{-p}+\delta^{1/3},
 \qquad p>1.
\]
This is a Dini modulus.  The classical Dini theorem makes its harmonic
conjugate continuous, hence its analytic projection lies in $\A$.  Finally,
the untruncated conjugate cutoff is anti-analytic, so its analytic projection
vanishes.  Splitting the circle into $E$ and its complement gives the desired
Cauchy transform supported on $E$.

\begin{theorem}[Log--log entropy construction]\label{thm:main}
Let $E\subsetneq\T$ be closed with $\m(E)>0$.  Suppose its complementary
arcs satisfy
\begin{equation}\label{eq:llentropy}
 \sum_n |I_n|\,\Phi(|I_n|)<\infty.
\end{equation}
Then there is an explicit bounded nonzero function $h$, supported on $E$,
such that $C_h\in\A$.
\end{theorem}

\section{The explicit formula}

Perform the Whitney decomposition from Lemma~2.1 of \cite{LM}:
\[
 I_n=\bigcup_{k\in\mathbb Z}B_{n,k},\qquad
 \ell_{n,k}:=|B_{n,k}|=\frac{|I_n|}{3\cdot2^{|k|}},
\]
where $\ell_{n,k}=\dist(B_{n,k},E)$.  Relabel these arcs as $B_j$, their
lengths as $\ell_j$, and their midpoints as $b_j\in\T$.  Put $r_j=1+\ell_j$.

For a fixed numerical $K>0$ chosen sufficiently large, define
\begin{equation}\label{eq:Fg}
 F_E(z)=-K\sum_j
 \frac{b_j\ell_j\Phi(\ell_j)}{r_jb_j-z},
 \qquad g_E(z)=\exp(F_E(z)),\qquad z\in\D.
\end{equation}
If $g_E^*$ denotes the almost-everywhere radial boundary value, set
\begin{equation}\label{eq:h}
 h_E(\zeta)=1_E(\zeta)\,
 \overline{\zeta g_E^*(\zeta)}.
\end{equation}
Equations \eqref{eq:Fg}--\eqref{eq:h} are the promised formula.

\section{Proof}

\begin{lemma}[Whitney summability]\label{lem:sum}
Condition \eqref{eq:llentropy} implies
\[
 \sum_j\ell_j\Phi(\ell_j)<\infty.
\]
\end{lemma}
\begin{proof}
For $L=|I_n|$, the contribution of its Whitney pieces is
\[
 \frac L3\Phi(L/3)+
 2\sum_{k\geq1}\frac{L}{3\cdot2^k}
 \Phi\!\left(\frac{L}{3\cdot2^k}\right).
\]
Since $\Phi(L/(3\cdot2^k))\leq C(\Phi(L)+\log(k+2))$ and both
$\sum2^{-k}$ and $\sum2^{-k}\log(k+2)$ converge, this is at most
$CL(1+\Phi(L))$.  Here $\Phi\geq1$, so summing in $n$ proves the claim.
\end{proof}

By Lemma~\ref{lem:sum}, the series in \eqref{eq:Fg} converges normally on
compact subsets of $\D$.  Moreover
\[
 \Re\frac{b_j}{r_jb_j-z}
 =\Re\frac1{r_j-\overline b_jz}>0,
\]
so $\Re F_E<0$ and $|g_E|\leq1$.  The function is nonzero in $\D$, and
the boundary uniqueness theorem for $H^\infty$ implies
$g_E^*\neq0$ almost everywhere on $\T$.  Consequently $h_E$ in
\eqref{eq:h} is bounded and nonzero.

The poles in \eqref{eq:Fg} cluster only at $E$, so $g_E$ continues
analytically across each compact subarc of $\T\setminus E$.  If
$\zeta\in B_j$, the $j$th summand and Whitney geometry give
\[
 \Re\frac{b_j}{r_jb_j-\zeta}\geq\frac{c}{\ell_j}.
\]
All summands have positive real part before multiplication by $-K$; hence
\begin{equation}\label{eq:decay}
 |g_E(\zeta)|\leq
 \exp\{-cK\Phi(\ell_j)\}
 =\bigl(e+\log(1/\ell_j)\bigr)^{-cK}.
\end{equation}
For $\zeta\in B_j$, $\ell_j\asymp\dist(\zeta,E)$.  Choosing $K$ so that
$p:=cK>2$, equation \eqref{eq:decay} becomes
\begin{equation}\label{eq:logdecay}
 |g_E(\zeta)|\leq
 C\bigl(\log(e/\dist(\zeta,E))\bigr)^{-p}.
\end{equation}

Define on the circle
\[
 S(\zeta)=1_{\T\setminus E}(\zeta)
 \overline{\zeta g_E^*(\zeta)},
\]
using the analytic continuation on $\T\setminus E$ and the value zero on
$E$.  We next verify the Dini modulus claimed above.  Write
$d(\zeta)=\dist(\zeta,E)$.  Whitney geometry also gives, uniformly in
$j$,
\[
 |r_jb_j-\zeta|\geq c\,d(\zeta),\qquad \zeta\in\T\setminus E.
\]
Indeed, a Whitney center close to $\zeta$ has length comparable with
$d(\zeta)$, while every other center is separated from $\zeta$ by a
constant multiple of $d(\zeta)$.  With
$M=\sum_j\ell_j\Phi(\ell_j)<\infty$, differentiation of
\eqref{eq:Fg} therefore yields
\begin{equation}\label{eq:deriv}
 |F_E'(\zeta)|\leq CKM d(\zeta)^{-2},
 \qquad |S'(\zeta)|\leq C_Ed(\zeta)^{-2}.
\end{equation}

Let two circle points be at arc distance at most $\delta$.  If either is
within $2\delta^{1/3}$ of $E$, then both are within $3\delta^{1/3}$ of
$E$, and \eqref{eq:logdecay} bounds their two values by
$C(\log(e/\delta))^{-p}$.  Otherwise the shorter connecting arc stays at
distance at least $\delta^{1/3}$ from $E$, and \eqref{eq:deriv} plus the
mean-value theorem bounds the difference by $C_E\delta^{1/3}$.  Thus
\begin{equation}\label{eq:modulus}
 \omega_S(\delta)\leq C_E\left[
 (\log(e/\delta))^{-p}+\delta^{1/3}\right].
\end{equation}
Since $p>1$, the right-hand side divided by $\delta$ is integrable near
zero.  Hence $S$ is Dini continuous.

The Dini theorem for conjugate functions states that the harmonic conjugate
of a Dini-continuous circle function is continuous \cite[Chapter III]{Zyg}.
Equivalently, its nonnegative-frequency Riesz projection $P_+S$ is the
boundary function of an element of $\A$.

Finally, $s:=\overline{\zeta g_E^*(\zeta)}$ belongs to
$\overline{zH^2}$, so $P_+s=0$.  Since $s=h_E+S$,
\[
 C_{h_E}=P_+h_E=-P_+S\in\A.
\]
If this function were zero, then $h_E\in\overline{zH^2}$; but $h_E$
vanishes on the nonempty open set $\T\setminus E$, so boundary uniqueness
would force $h_E=0$, a contradiction.  This completes the proof of
Theorem~\ref{thm:main}. \qed

\section{Strict extension beyond Beurling--Carleson sets}

\begin{example}\label{ex:strict}
For each sufficiently large integer $k$, take
$N_k=\lfloor a e^k/k^2\rfloor$ disjoint arcs of length $e^{-k}$, with
$a>0$ small enough that their total length is less than one, and let $E$ be
the complement of their union.  Then $\m(E)>0$ and
\[
 \sum_kN_ke^{-k}\Phi(e^{-k})
 \asymp\sum_k\frac{\log(k+e)}{k^2}<\infty,
\]
whereas
\[
 \sum_kN_ke^{-k}\log(e^k)
 \asymp\sum_k\frac1k=\infty.
\]
Thus $E$ satisfies Theorem~\ref{thm:main} but is not a
Beurling--Carleson set.  Hence the explicit formula \eqref{eq:Fg} applies
directly to closed sets outside the hypothesis of the source paper's
cutoff construction.
\end{example}

\section{Verification, novelty, and limitations}

The proof uses only positive-term estimates.  The accompanying script checks
the convergence/divergence separation in Example~\ref{ex:strict} and the
bounded ratio between the Whitney and original log--log entropy sums.  It is
a consistency check, not part of the proof.

A bounded literature search through 27 August 2026 used arXiv:2111.14112 and
the phrases ``explicit Cauchy transform closed set'', ``Dini Cauchy transform
Beurling--Carleson'', ``log log Beurling--Carleson'', and ``loglog entropy
disk algebra''.  It found the source article and standard Dini-regular
Cauchy-transform results, but no matching weakening of the source's entropy
hypothesis.  Novelty confidence is moderate, not definitive.

The theorem does not solve the stated problem for every positive-measure
closed set.  The remaining obstruction is exactly the possible divergence
of \eqref{eq:llentropy}; the present positive exterior-pole construction then
cannot simultaneously have finite pole mass and a uniform Dini modulus.

\paragraph{Human-review recommendation.}
Review as a likely-valid substantial partial result.  The highest-value
checks are the denominator estimate preceding \eqref{eq:deriv}, the two-case
Dini modulus proof, and the boundary uniqueness step proving nontriviality.

\begin{thebibliography}{9}
\bibitem{LM}
A.~Limani and B.~Malman,
\emph{Constructions of some families of smooth Cauchy transforms},
Canadian Journal of Mathematics \textbf{76} (2024), 319--344;
arXiv:2111.14112.

\bibitem{Zyg}
A.~Zygmund,
\emph{Trigonometric Series}, Vol.~I,
Cambridge University Press, third edition, 2002.
\end{thebibliography}

\end{document}
